\section{Conclusion}

The experiments answer the two study questions in different ways. For Q1, 
different combinations of read concern, write concern, and causal
sessions produced different client-visible behaviour. Under controlled replica
divergence, C1 (\texttt{local}, \texttt{w:1}, no causal consistency) produced
counterexamples for all four client-centric properties. In contrast, C6
(\texttt{majority}, \texttt{majority}, causal consistency) produced no
decidable adversarial histories in the configuration campaign: the RYW and MR
reads timed out, while the MW and WFR histories had unresolved write effects.
The intermediate configurations further showed that read concern, write
concern, and causal sessions address different parts of the client's operation
history.

For Q2, the type of failure was important. Secondary loss (F1) and operations
performed after a primary election (F2) produced no violations in the tested
schedules. All observed fault-campaign violations occurred under the
replication-path partition (F3), which produced 56 violations in the tested
histories. The partition separated the isolated former primary from the
majority side, allowing the two sides to expose different states. This made
the effects of the tested consistency settings most visible.

Overall, the experiments show why client-centric consistency cannot be
evaluated only by counting observed violations. Stronger settings may prevent
an inconsistent value from being returned by delaying an operation or leaving
the history unresolved instead. Consistency results should therefore be
interpreted together with operation completion, latency, and whether the
history is decidable. These conclusions are limited to the tested schedules,
timeout settings, and three-member single-host Docker deployment.